\documentclass[11pt,a4paper]{article}

% --- Paquetes Fundamentales ---
\usepackage[utf8]{inputenc}
\usepackage[T1]{fontenc}
\usepackage[spanish,es-tabla]{babel}
\usepackage[margin=2.5cm, headheight=15pt]{geometry} 
\usepackage{graphicx}
\usepackage{fancyhdr}
\usepackage{xcolor}
\usepackage[colorlinks=true, linkcolor=blue, urlcolor=blue]{hyperref}

% --- Matemáticas y Electrónica ---
\usepackage{amsmath, amssymb}
\usepackage{siunitx}
\usepackage[american]{circuitikz}

% --- Elementos de Diseño y Código ---
\usepackage{tcolorbox}
\usepackage{booktabs}
\usepackage{enumitem}
\usepackage{listings}

% --- Configuración de siunitx ---
\sisetup{
    output-decimal-marker = {,},
    per-mode = symbol,
    inter-unit-product = \ensuremath{{}\cdot{}}
}

% --- Configuración de Python (Listings) ---
\lstset{
    language=Python,
    basicstyle=\ttfamily\footnotesize,
    keywordstyle=\color{blue}\bfseries,
    stringstyle=\color{teal},
    commentstyle=\color{green!50!black}\itshape,
    backgroundcolor=\color{gray!5},
    frame=single,
    rulecolor=\color{gray!40},
    breaklines=true
}

% --- Cajas Personalizadas para Dinámica Docente ---
\newtcolorbox{acciondocente}[1][]{
    colback=red!5!white,
    colframe=red!75!black,
    fonttitle=\bfseries,
    title=Acción Docente,
    #1
}

\newtcolorbox{justificacion}[1][]{
    colback=blue!5!white,
    colframe=blue!75!black,
    fonttitle=\bfseries,
    title=Justificación Pedagógica (CDIO),
    #1
}

% --- Estilo de Página ---
\pagestyle{fancy}
\fancyhf{}
\lhead{\textbf{Introducción a la Mecatrónica}}
\rhead{Ingeniería Mecatrónica -- UCB Tarija}
\cfoot{Página \thepage}
\renewcommand{\headrulewidth}{0.4pt}

% --- Metadatos del Documento ---
\title{\vspace{-1.5cm}\Large\textbf{Guía Docente Paso a Paso: Sesión 03}\\
\large Redes de Impedancia Compleja y Modelado de Transductores Térmicos}
\author{M.Sc. Ing. Bernardo Quiroga Turdera}
\date{Lunes, 10 de agosto de 2026 \quad|\quad 09:00 - 11:00 (120 min)}

\begin{document}

\maketitle
\vspace{-1cm}

\begin{justificacion}
Esta clase teórica es la piedra angular del Módulo 1. Prepara de forma directa a los estudiantes para el Laboratorio 1 derivando matemáticamente el álgebra de impedancias complejas y demostrando cómo los efectos parásitos alteran la lectura del sensor.
\end{justificacion}

\section*{I. Agenda y Distribución del Tiempo (120 Minutos)}
\begin{itemize}[label=$\circ$, itemsep=2pt]
    \item \textbf{09:00 - 09:15 (15 min):} Retroalimentación de la Sesión 2 y puente al dominio fasorial.
    \item \textbf{09:15 - 09:55 (40 min):} Pizarra 1 - Álgebra de Impedancias y Leyes de Kirchhoff en CA.
    \item \textbf{09:55 - 10:35 (40 min):} Pizarra 2 - Modelado del Sensor NTC y Parásitos.
    \item \textbf{10:35 - 10:50 (15 min):} Taller de Pizarra Guiado - Cálculo analítico a \qty{30}{\celsius}.
    \item \textbf{10:50 - 11:00 (10 min):} Cierre, entrega de Script base (Python) y roles de laboratorio.
\end{itemize}

\vspace{0.5cm}
\hrule
\vspace{0.5cm}

\section*{II. Desarrollo Teórico: Guía de Pizarra (09:15 - 10:35)}

\subsection*{Pizarra 1: Álgebra de Impedancias y Leyes de Kirchhoff (09:15 - 09:55)}
La relación entre la tensión $\hat{V}$ y corriente $\hat{I}$ fasoriales generaliza la resistencia hacia la Impedancia Compleja ($Z$):
\begin{equation}
    \hat{V} = Z \cdot \hat{I} \implies Z = \frac{\hat{V}}{\hat{I}} = \vert{}Z\vert{} \angle \theta = R + jX
\end{equation}
\begin{itemize}
    \item \textbf{Resistencia ($R$):} Parte real (disipa energía en forma de calor).
    \item \textbf{Reactancia ($X$):} Parte imaginaria (almacena energía electromagnética).
    \begin{itemize}
        \item Inductiva ($X > 0$): $X_L = \omega L$
        \item Capacitiva ($X < 0$): $X_C = -\frac{1}{\omega C}$
    \end{itemize}
\end{itemize}

Leyes de Kirchhoff y Divisores de Tensión se aplican idénticamente, sumando fasores:
\begin{equation}
    \sum \hat{V}_k = 0 \quad \text{y} \quad \hat{V}_{Z2} = \hat{V}_{in} \left( \frac{Z_2}{Z_1 + Z_2} \right)
\end{equation}

\subsection*{Pizarra 2: Modelado del Sensor NTC y Efectos Parásitos (09:55 - 10:35)}

\begin{center}
    \begin{circuitikz}[scale=1.0, transform shape]
        \draw (0,4) to[sV, l=$V_{in}$ (CA)] (0,0);
        \draw (0,4) -- (2,4) -- (5,4);
        \draw (0,0) -- (2,0) -- (5,0);
        \draw (2,0) node[ground]{};
        
        \draw (2,4) to[R, l=$R_1$] (2,2) to[R, l=$R_2$] (2,0);
        \draw (5,4) to[R, l=$R_3$] (5,2) to[generic, l=$Z_{sensor}$] (5,0);
        
        \draw (2,2) to[short, *-o] (3,2) node[above]{$+$};
        \draw (5,2) to[short, *-o] (4,2) node[above]{$-$};
        \node at (3.5, 2) {$\hat{V}_{out}$};
    \end{circuitikz}
\end{center}

Ecuación Steinhart-Hart reducida (Comportamiento del NTC):
\begin{equation}
    R_{NTC}(T) = R_0 \cdot e^{\beta \left( \frac{1}{T} - \frac{1}{T_0} \right)}
\end{equation}

Modelado de la línea de transmisión (15m trenzado, con $R_{cable}$, $L_{cable}$, $C_{para}$):
\begin{equation}
    Z_{sensor}(j\omega) = R_{cable} + j\omega L_{cable} + \left( R_{NTC}(T) \parallel \frac{1}{j\omega C_{para}} \right)
\end{equation}

Tensión de Desbalance Fasorial del Puente:
\begin{equation}
    \hat{V}_{out}(j\omega) = \hat{V}_{in} \left( \frac{Z_{sensor}(j\omega)}{R_3 + Z_{sensor}(j\omega)} - \frac{R_2}{R_1 + R_2} \right)
\end{equation}

\newpage
\section*{III. Taller de Pizarra Guiado: Resolución Paso a Paso (10:35 - 10:50)}

\begin{acciondocente}
Divide la pizarra en secuencias numéricas. Pide a los estudiantes que configuren sus calculadoras en modo complejo (grados) y vayan dictándote los resultados de cada etapa.
\end{acciondocente}

\noindent \textbf{Datos iniciales del problema planteado:}
\begin{itemize}[label=$\circ$, itemsep=1pt]
    \item Excitación: $V_{in} = \qty{2}{\volt}_{pp} \implies \hat{V}_{in} = 1 \angle 0^\circ \text{ V}$ (pico). Frecuencia: $f_c = \qty{10}{\kilo\hertz}$ ($\omega_c \approx \qty{62831.85}{\radian\per\second}$).
    \item Resistencias del puente: $R_1 = R_2 = R_3 = \qty{10}{\kilo\ohm}$.
    \item Condición física: Cámara a $T = \qty{30}{\celsius}$ (\qty{303.15}{\kelvin}).
    \item Parámetros de diseño: $R_0 = \qty{10}{\kilo\ohm}$ a \qty{298.15}{\kelvin}, $\beta = \qty{3950}{\kelvin}$.
\end{itemize}

\vspace{0.3cm}
\noindent \textbf{Paso 1: Cálculo de la resistencia del NTC a \qty{30}{\celsius}}
\begin{equation*}
    R_{NTC}(303.15) = 10000 \cdot e^{3950 \left( \frac{1}{303.15} - \frac{1}{298.15} \right)} \approx \qty{8059.2}{\ohm}
\end{equation*}

\noindent \textbf{Paso 2: Cálculo de Reactancias a \qty{10}{\kilo\hertz}}
\begin{align*}
    X_L &= \omega_c L_{cable} = 62831.85 \cdot 15 \cdot 10^{-6} \approx \qty{0.94}{\ohm} \quad \text{\small(Despreciable frente a los k$\Omega$)} \\
    Z_C &= \frac{1}{j\omega_c C_{para}} = \frac{1}{j(62831.85)(4.7 \cdot 10^{-9})} = -j 3386.2\ \Omega
\end{align*}

\noindent \textbf{Paso 3: Cálculo del Paralelo ($R_{NTC} \parallel Z_C$)}
\begin{equation*}
    Z_{paralelo} = \frac{8059.2 \cdot (-j 3386.2)}{8059.2 - j 3386.2} \approx 1208.6 - j 2876.5\ \Omega
\end{equation*}

\noindent \textbf{Paso 4: Impedancia Total del Sensor ($Z_{sensor}$)}
Agregando la resistencia parásita del cable ($R_{cable} = \qty{10}{\ohm}$):
\begin{equation*}
    Z_{sensor} \approx 1218.6 - j 2875.56\ \Omega = 3123.8 \angle -67.02^\circ\ \Omega
\end{equation*}

\noindent \textbf{Paso 5: Tensión Fasorial de Salida ($\hat{V}_{out}$)}
Sustituyendo en la ecuación del puente:
\begin{equation*}
    \hat{V}_{out} = 1 \angle 0^\circ \left( \frac{3123.8 \angle -67.02^\circ}{10000 + 1218.6 - j 2875.56} - 0.5 \right) \approx 0.252 \angle -162.4^\circ\text{ V}
\end{equation*}

\begin{justificacion}
\textbf{Conclusión de la clase:} Destaca a los estudiantes que la salida no solo cambia en amplitud cuando varía la temperatura, sino que presenta un severo \textbf{desfase angular negativo ($\theta \approx -162.4^\circ$)} causado por la capacitancia parásita del cable. ¡Este es el fenómeno exacto que medirán en el osciloscopio el miércoles!
\end{justificacion}

\section*{IV. Cierre y Tarea para el Miércoles (10:50 - 11:00)}

\noindent \textbf{Código Base de Python para el Laboratorio 1:}
\begin{lstlisting}
import numpy as np
import pandas as pd

# Carga de traza CSV exportada del osciloscopio
df = pd.read_csv("lab1_medicion_30C.csv")

t = df['Time'].values
v_in = df['Channel1'].values
v_out = df['Channel2'].values

# Calculo de amplitud pico experimental
v_out_pico = np.max(v_out)

# Medicion automatica de diferencia de tiempo entre picos (Delta t)
idx_pico_in = np.argmax(v_in)
idx_pico_out = np.argmax(v_out)
delta_t = t[idx_pico_out] - t[idx_pico_in]

f_c = 10000.0  # Frecuencia portadora 10 kHz
fase_grados = delta_t * f_c * 360.0

print(f"Amplitud V_out experimental: {v_out_pico:.3f} V")
print(f"Desfase experimental: {fase_grados:.2f} grados")
\end{lstlisting}

\vspace{0.3cm}
\begin{tcolorbox}[colback=gray!5, colframe=black!50, title=\textbf{Lista de Verificación de Hardware para el Miércoles (Lab 1)}]
\small
Recordar a los grupos traer su kit completo: 1x Termistor NTC (\qty{10}{\kilo\ohm} a \qty{25}{\celsius}), 3x Resistores \qty{10}{\kilo\ohm} al 1\%, 1x Capacitor de poliéster de \qty{4.7}{\nano\farad}, Protoboard, cables, y memoria USB.
\end{tcolorbox}

\end{document}